\documentclass[10pt,a4paper]{amsart}
\usepackage[margin=19mm]{geometry}
\usepackage{amsmath,amssymb,mathtools}
\usepackage[scr=rsfs,cal=euler]{mathalfa}
\usepackage{xfrac}
\usepackage[unicode,hidelinks,bookmarks=false]{hyperref}
\newcommand{\ie}{i.e.\ }
\newcommand{\cf}{cf.\ }
\newcommand{\resp}{respectively\ }
\newcommand{\Subsection}{\subsection}
\providecommand{\nobreakdash}{\nobreak\hbox{-}}
\setlength{\emergencystretch}{2em}
\multlinegap0pt
\usepackage{amssymb}
\usepackage{enumerate}
\usepackage[normalem]{ulem}
\usepackage{bm}
\usepackage{float}
\usepackage{mathtools}
\usepackage{tikz}
\newtheorem{proposition}{Proposition}
\usepackage{cleveref}
\crefname{proposition}{Proposition}{Propositions}
\newcommand{\bR}{\mathbb{R}}
\def\cR{{\mathscr R}}
\newcommand\scalemath[2]{\scalebox{#1}{\mbox{\ensuremath{\displaystyle #2}}}}
\def\ua{\uparrow}
\def\wh{\widehat}
\def\wt{\widetilde}
\title{Estimating the roughness exponent of stochastic volatility from discrete observations of the integrated variance}
\author{Xiyue Han$^*$ and Alexander Schied}
\date{Source: arXiv:2307.02582v4 (CC BY 4.0)}
\begin{document}
\maketitle

\noindent\emph{Source and licence.} Estimating the roughness exponent of stochastic volatility from discrete observations of the integrated variance. Version arXiv:2307.02582v4 (CC BY 4.0), distributed by arXiv under the Creative Commons Attribution 4.0 International licence, http://creativecommons.org/licenses/by/4.0/, as recorded in the arXiv licence metadata for this exact version and shown on the abstract page https://arxiv.org/abs/2307.02582v4. Full licence notice: \texttt{licenses/arxiv-2307.02582-CC-BY-4.0.txt}.\par\smallskip

\noindent\textit{Editorial scope.} Counted unit: the article's proposition z (source0.tex lines 1023--1033). The article's complete original proof of it is delivered with it (source0.tex lines 1101--1194, one \texttt{proof} environment of 94 lines, which the article places away from the statement). No mathematics is written, repaired or re-derived: the statement and the proof are the article's own lines, in the article's own notation, and no retained line is substituted or re-flowed.  Retained uncounted as the source-local context the proof needs: lines 1014--1019.  Nothing was omitted from any retained span, so the package carries no omission record.  No retained line contains a citation, so the wrapper carries no bibliography and no bibliography entry.  No retained line contains a figure environment, an image-inclusion command or drawing code, so no image is delivered and the package declares no asset dependency.  Editorial formatting only, in addition: an AMS \texttt{amsart} preamble that restates the article's own declarations and macros used by the retained lines, namely the environments \texttt{proposition} on separate counters, together with the standard packages those lines need.  The retained environments print in this document as Proposition 1 in order of appearance, whereas the article prints the same items with the numbering of its own full document.  The complete native source of the article as distributed in the arXiv e-print for this exact version is bundled unchanged as \texttt{sources/143896/source0.tex}.
\begin{equation}\label{coefficients depending on f eq}
	\begin{split}
		\theta^f_{n,k}& = 2^{n/2}\left(2f\Big(\frac{2k+1}{2^{n+1}}\Big)-f\Big(\frac{k}{2^{n}}\Big) - f\Big(\frac{k+1}{2^{n}}\Big)\right), \\
		\vartheta^f_{n,k}&= 2^{3n/2+3}\left(f\Big(\frac{4k}{2^{n+2}}\Big)-2f\Big(\frac{4k+1}{2^{n+2}}\Big)+2f\Big(\frac{4k+3}{2^{n+2}}\Big)-f\Big(\frac{4k+4}{2^{n+2}}\Big)\right).
	\end{split}
\end{equation}

\begin{proposition}\label{Prop tilde H z} Suppose there exists ${R}\in(0,1)$ such that the following conditions hold.
\begin{enumerate}
\item \label{Prop tilde H z part a}We have  
	\begin{equation}\label{Prop tilde H z a assumption eq}
	0< \liminf_{n \ua \infty}2^{n(2{R}-2)}\sum_{k = 0}^{2^n-1}\left(\vartheta^y_{n,k}\right)^2\le \limsup_{n \ua \infty}2^{n(2{R}-2)}\sum_{k = 0}^{2^{n}-1}\left(\vartheta^y_{n,k}\right)^2<\infty.
	\end{equation}
	
\item\label{Prop tilde H z part b}  The function $x$ is H\"older continuous with exponent $\alpha\in(2{R}/5,1]$. 
\end{enumerate}
Then, if $g\in C^2(\bR)$ is strictly monotone,  we have $\lim_n\wh\cR_n(v)={R}$.
\end{proposition}

\begin{proof}[Proof of \Cref{Prop tilde H z}]
	In this proof, we will work with the actual and estimated Faber--Schauder coefficients of the functions $x$,  $u=g\circ x$, and their respective antiderivatives $y$ and $v$. Here, we will use a superscript as in \eqref{coefficients depending on f eq} to indicate from which function the Faber--Schauder coefficients will be computed. Our goal in this proof is to show that \eqref{Prop tilde H z a assumption eq} carries over to the coefficients $\vartheta^v_{n,k}$.
	That is, 
	\begin{equation}\label{Prop tilde H z first goal eq}
		0<\liminf_{n \ua \infty}2^{n(2{R}-2)}\sum_{k = 0}^{2^{n}-1}\big(\vartheta^v_{n,k}\big)^2\le \limsup_{n \ua \infty}2^{n(2{R}-2)}\sum_{k = 0}^{2^{n}-1}\big(\vartheta^v_{n,k}\big)^2<\infty.	\end{equation}
Once this has been achieved, taking logarithms in \eqref{Prop tilde H z first goal eq}, dividing by $2n$, and passing to the limit will then yield
	${R}-\wh\cR_n(v)\to0$, which is the assertion. 
	
	It remains to establish \eqref{Prop tilde H z first goal eq}. 
	Rewriting the second line in \eqref{coefficients depending on f eq}
	gives after a short computation that 
	\begin{equation}\label{eq FS y}
		\vartheta^f_{n,k} =   2^{n+5/2}\left(\theta^f_{n+1,2k+1}-\theta^f_{n+1,2k}\right).	\end{equation}
	Let us  introduce the notation $\theta_{m,k}^f(s):=\theta_{m.k}^{f(s+\cdot)}$. That is, $\theta_{m,k}^f(s)$ are the Faber--Schauder coefficients of the function $t\mapsto f(s+t)$ for given $s\ge0$. Here and in the sequel, we will avoid undefined arguments of functions in case $s+t>1$ by setting $f(t):= f(t \wedge 1)$ for $t \ge 1$. With this notation, 
	we get from \eqref{eq FS y} that for $f\in C^1[0,\infty)$,		\begin{equation}\label{eq zeta v}
		\vartheta^f_{n,k} = 2^{n+5/2}\int_{0}^{2^{-n-1}}\theta^{f'}_{n+1,2k}(s)\,ds .
	\end{equation}
	Applying the mean-value theorem and the intermediate value theorem yields certain intermediate times $\tau_{n+2,k}(s) \in [2^{-n-2}k+s, 2^{-n-2}(k+1)+s]$ such  that for $s \in [0,2^{-n-1}]$, 
	\begin{equation*}
		\begin{split}
			\theta^{u}_{n+1,2k}(s) &= 2^{(n+1)/2}\left(2u\Big(\frac{4k+1}{2^{n+2}}+s\Big)-u\Big(\frac{4k}{2^{n+2}}+s\Big)-u\Big(\frac{4k+2}{2^{n+2}}+s\Big)\right)\\&= 2^{(n+1)/2} g'\big(x(\tau_{n+2,4k}(s))\big)\left(x\Big(\frac{4k+1}{2^{n+2}}+s\Big)-x\Big(\frac{4k}{2^{n+2}}+s\Big)\right) \\
			&+2^{(n+1)/2} g'\big(x(\tau_{n+2,4k+1}(s))\big)\left(x\Big(\frac{4k+1}{2^{n+2}}+s\Big)-x\Big(\frac{4k+2}{2^{n+2}}+s\Big)\right),\\
			&=2^{(n+1)/2}\scalemath{0.9}{\left(\frac{g'\big(x(\tau_{n+2,4k}(s))\big)+g'\big(x(\tau_{n+2,4k+1}(s))\big)}{2}\left(2x\Big(\frac{4k+1}{2^{n+2}}+s\Big)-x\Big(\frac{4k}{2^{n+2}}+s\Big)-x\Big(\frac{4k+2}{2^{n+2}}+s\Big)\right)\right)}\\&\quad+ 2^{(n+1)/2}\left(\frac{g'\big(x(\tau_{n+2,4k}(s))\big)-g'\big(x(\tau_{n+2,4k+1}(s))\big)}{2}\left(x\Big(\frac{4k+2}{2^{n+2}}+s\Big)-x\Big(\frac{4k}{2^{n+2}}+s\Big)\right)\right).
		\end{split}
	\end{equation*}
	The intermediate value theorem  and the mean-value theorem  also imply that there are intermediate times $\tau^\sharp_{n+1,2k}(s),\tau^\flat_{n+1,2k}(s) \in [2^{-n-1}2k+s,2^{-n-1}(2k+1)+s]$ such that 
	\begin{align*}
		\frac{1}{2}\left(g'\big(x(\tau_{n+2,4k}(s))\big)+g'\big(x(\tau_{n+2,4k+1}(s))\big)\right) &= g'\big(\tau^\sharp_{n+1,2k}(s)\big),\\
		\frac{1}{2}\left(g'\big(x(\tau_{n+2,4k}(s))\big)-g'\big(x(\tau_{n+2,4k+1}(s))\big)\right) &= \frac{1}{2}g''\big(\tau^\flat_{n+1,2k}(s)\big)\left(x(\tau_{n+2,4k}(s))-x(\tau_{n+2,4k+1}(s))\right).
	\end{align*} 
	With the shorthand notation
	\begin{equation*}
		\zeta^x_{n+1,2k}(s):= 2^{(n+1)/2}\left(x\Big(\frac{4k+2}{2^{n+2}}+s\Big)-x\Big(\frac{4k}{2^{n+2}}+s\Big)\right)\big(x(\tau_{n+2,4k}(s))-x(\tau_{n+2,4k+1}(s))\big),
	\end{equation*}
	we then have
	\begin{equation*}
		\theta^{u}_{n+1,2k}(s) =g'\big(\tau^\sharp_{n+1,2k}(s)\big)\theta^{x}_{n+1,2k}(s) + g''\big(\tau^\flat_{n+1,2k}(s)\big)\zeta^{x}_{n+1,2k}(s).
	\end{equation*}
	Plugging the preceding equation into \eqref{eq zeta v} and applying the mean value theorem for integrals yields intermediate times $\tau^\sharp_{n+1,k}, \tau^\flat_{n+1,k} \in [2^{-n-1}{k},2^{-n-1}(k+1)]$  that are independent of $s$ such that
	\begin{equation*}
		\begin{split}
			\vartheta^v_{n,k} &=  2^{n+5/2}g'\big(x(\tau^\sharp_{n+1,2k})\big)\int_{0}^{2^{-n-1}}\theta^{x}_{n+1,2k}(s)\,ds + 2^{n+5/2}g''\big(x(\tau^\flat_{n+1,2k})\big)\int_{0}^{2^{-n-1}} \zeta^{x}_{n+1,2k}(s)\,ds \\&= g'\big(x(\tau^\sharp_{n+1,2k})\big)\vartheta^y_{n,k} + 2^{n+5/2}g''\big(x(\tau^\flat_{n+1,2k})\big)\int_{0}^{2^{-n-1}} \zeta^{x}_{n+1,2k}(s)\,ds.
		\end{split}
	\end{equation*}
	Introducing the shorthand notation	\begin{equation*}
		\wt \zeta^x_{n+1,2k}:= 2^{n+5/2}\int_{0}^{2^{-n-1}} \zeta^{x}_{n+1,2k}(s)\,ds,
	\end{equation*}
	we can write 
	\begin{equation}\label{three terms eq}
		\begin{split}
			\big(	\vartheta^v_{n,k} \big)^2
			& = \big(g'\big(x(\tau^\sharp_{n+1,2k})\big)\big)^2\big(\vartheta^y_{n,k}\big)^2 + \left(g''\big(x(\tau^\flat_{n+1,2k})\big)\right)^2\big(\wt\zeta^x_{n+1,2k}\big)^2\\
			&\qquad + 2g'\big(x(\tau^\sharp_{n+1,2k})\big)g''\big(x(\tau^\flat_{n+1,2k})\big)\vartheta^y_{n,k}\wt\zeta^x_{n+1,2k}.	\end{split}
	\end{equation}
	For each of the three terms on the right,  we will now analyze its contribution to the quantities in~\eqref{Prop tilde H z first goal eq}. 
	The main contribution comes from the first term on the right.	Indeed, our assumptions on $g$ imply that there are constants $0<c_-\le c_+<\infty$ such that $c_-<(g'(x(t)))^2<c^+$ for all $t\in[0,1]$, and so 	
	$$c_-2^{n(2{R}-2)}\sum_{k = 0}^{2^{n}-1}\big(\vartheta^y_{n,k}\big)^2\le 2^{n(2{R}-2)}\sum_{k = 0}^{2^{n}-1}\big(\vartheta^v_{n,k}\big)^2\le c_+2^{n(2{R}-2)}\sum_{k = 0}^{2^{n}-1}\big(\vartheta^y_{n,k}\big)^2.
	$$
	This will establish \eqref{Prop tilde H z first goal eq} as soon as
	we have shown that the contributions of the two remaining terms in \eqref{three terms eq} are asymptotically negligible.	For the second term, we use the H\"older continuity of $x$ to get a constant $c_x$ for which
	\begin{equation*}
		|x(\tau_{n+2,4k}(s))-x(\tau_{n+2,4k+1}(s))| \le c_x|\tau_{n+2,4k}(s) - \tau_{n+2,4k+1}(s)|^\alpha \le c_x2^{-\alpha n}
	\end{equation*}
	Furthermore,  there exists $\kappa_x > 0$ such that $32 (g''(x(s)))^2 \le \kappa_x$ for all $s\in[0,1]$. Thus,
	\begin{equation*}
		\begin{split}
			&2^{(2{R}-2)n}\sum_{k = 0}^{2^{n}-1}\left(g''\big(x(\tau^\flat_{n+1,2k})\big)\right)^2\left(\wt\zeta_{n+1,2k}
			\right)^2 \\
			&= 2^{(2{R}-2)n}\sum_{k = 0}^{2^{n}-1}\left(g''\big(x(\tau^\flat_{n+1,2k})\big)\right)^2\left(2^{n+5/2}\int_{0}^{2^{-n-1}}\zeta^{x}_{n+1,2k}(s)\,ds\right)^2\\
			&\le \kappa_x 2^{2{R}n}\sum_{k = 0}^{2^{n}-1}\left(\int_{0}^{2^{-n-1}}\zeta^{x}_{n+1,2k}(s)\,ds\right)^2 \le \kappa_x 2^{2{R}n}\sum_{k = 0}^{2^{n}-1}2^{-n-1}\int_{0}^{2^{-n-1}}\left(\zeta^{x}_{n+1,2k}(s)\right)^2\,ds \\
			&\le \kappa_x 2^{(2{R}-1)n}\int_{0}^{2^{-n-1}}\sum_{k = 0}^{2^{n}-1}\left(x\Big(\frac{4k+2}{2^{n+2}}+s\Big)-x\Big(\frac{4k}{2^{n+2}}+s\Big)\right)^2\left(x(\tau_{n+2,4k}(s))-x(\tau_{n+2,4k+1}(s))\right)^2\,ds \\
			&= \kappa_x 2^{\alpha n}\int_{0}^{2^{-n-1}}2^{(2{R}-1-\alpha)n}\sum_{k = 0}^{2^{n}-1}\left(x\Big(\frac{2k+1}{2^{n+1}}+s\Big)-x\Big(\frac{2k}{2^{n+1}}+s\Big)\right)^2\big(c^2_x2^{-2\alpha n}\big)\,ds \\
			&\le \kappa_x c^2_x \int_{0}^{2^{-n-1}}\sup_{s \in [0,2^{-n-1}]}\left(2^{(2{R}-1-3\alpha)n}\sum_{k = 0}^{2^{n+1}-1}\left(x\Big(\frac{2k+1}{2^{n+1}}+s\Big)-x\Big(\frac{2k}{2^{n}}+s\Big)\right)^2\right)\,ds.	\end{split}
	\end{equation*}
	Moreover, \ref{Prop tilde H z part b} implies the  integrand in the final term converges to zero:
	\begin{equation}\label{Prop tilde H z aux claim}
		\lim_{n\ua\infty}\sup_{0\le s\le 2^{-n-1}}2^{(2{R}-1-3\alpha)n}\sum_{k = 0}^{2^{n+1}-1}\left(x\Big(\frac{2k+1}{2^{n+1}}+s\Big)-x\Big(\frac{2k}{2^{n}}+s\Big)\right)^2=0.
	\end{equation}
	Indeed, by the H\"older continuity of $x$, we can again use the  constant $c_x$ to get 
	\begin{align*}
		2^{(2{R}-1-3\alpha)n}\sum_{k = 0}^{2^{n+1}-1}\left(x\Big(\frac{2k+1}{2^{n+1}}+s\Big)-x\Big(\frac{2k}{2^{n}}+s\Big)\right)^2&\le 2^{(2{R}-1-3\alpha)n}\cdot 2^{n+1}\cdot c_x^22^{-2(n+1)\alpha};
	\end{align*}
	the right-hand side is equal to $c_x^2\cdot2^{1-2\alpha}\cdot2^{(2{R}-5\alpha)n}$, which converges to zero as $n\ua\infty$. Altogether, this shows that the contribution of the second term on the right-hand side of \eqref{three terms eq} is negligible. 
	
	
	
	For the cross-product term on the rightmost side of \eqref{three terms eq}, we get from the Cauchy--Schwarz inequality, 
	\begin{equation*}
		\begin{split}
			\lefteqn{\lim_{n \ua \infty}2^{(2{R}-2)n}\sum_{k = 0}^{2^n-1}g'\big(x(\tau^\sharp_{n+1,2k})\big)g''\big(x(\tau^\flat_{n+1,2k})\big)\vartheta_{n,k}\wt\zeta^x_{n+1,2k}}
			\\&\le \sqrt{\lim_{n \ua \infty}2^{(2{R}-2)n}\sum_{k = 0}^{2^n-1}\left(g'\big(x(\tau^\sharp_{n+1,2k})\big)\right)^2\big(\vartheta_{n,k}\big)^2} \sqrt{\lim_{n \ua \infty}2^{(2{R}-2)n}\sum_{k = 0}^{2^n-1}\left(g''\big(x(\tau^\flat_{n+1,2k})\big)\right)^2\big(\wt\zeta^x_{n+1,2k}\big)^2} = 0.
		\end{split}
	\end{equation*}
	Altogether, \eqref{Prop tilde H z first goal eq} follows.\end{proof}

\end{document}
